\begin{ex}
	Thí nghiệm của Brown quan sát
	\choice
	{hạt phấn hoa bằng kính hiển vi}
	{\True chuyển động của hạt phấn hoa trong nước bằng kính hiển vi}
	{cánh hoa trong nước bằng kính hiển vi}
	{chuyển động của cánh hoa}
	\loigiai{
		\begin{itemize}
			\item Thí nghiệm của Brown quan sát chuyển động của các hạt phấn hoa trong nước, thấy các hạt phấn chuyển động hỗn độn, không ngừng
		\end{itemize}
	}
	\haidong{3}
\end{ex}
\begin{ex}
	Trong thí nghiệm của Brown, các hạt phấn hoa chuyển động hỗn độn không ngừng vì
	\choice
	{giữa chúng có khoảng cách}
	{chúng là các phân tử}
	{\True các phân tử nước chuyển động không ngừng, va chạm vào chúng từ mọi phía}
	{chúng là các thực thể sống}
	\loigiai{
		\begin{itemize}
			\item Trong thí nghiệm của Brown, sở dĩ các hạt phấn hoa lại chuyển động hỗn độn không ngừng là vì các phân tử nước chuyển động không ngừng va chạm vào các hạt phấn hoa từ mọi phía.
		\end{itemize}
	}
	\haidong{3}
\end{ex}

\begin{ex}
	Phát biểu nào sau đây là \textbf{sai} khi nói về cấu tạo chất?	
	\choice
	{Các phân tử luôn luôn chuyển động không ngừng}
	{Các phân tử chuyển động càng nhanh thì nhiệt độ của vật càng cao}
	{\True Các phân tử luôn luôn đứng yên và chỉ chuyển động khi nhiệt độ của vật cao}
	{Các chất được cấu tạo từ các hạt riêng biệt gọi là phân tử}
	\loigiai{
		\begin{itemize}
			\item Các phân tử chuyển động không ngừng.
			Chuyển động của các phân tử được gọi là chuyển động nhiệt.
		\end{itemize}
	}
	\haidong{2}
\end{ex}

\begin{ex}
	Nguyên tử, phân tử \textbf{không} có tính chất nào sau đây?
	\choice
	{Chuyển động không ngừng}
	{Giữa chúng có khoảng cách}
	{\True Nở ra khi nhiệt độ tăng, co lại khi nhiệt độ giảm}
	{Chuyển động càng nhanh khi nhiệt độ càng cao}
	\loigiai{
		\begin{itemize}
			\item Nguyên tử, phân tử không có tính chất nở ra khi nhiệt độ tăng, co lại khi nhiệt độ giảm, chỉ có chuyển động càng nhanh khi nhiệt độ càng cao
		\end{itemize}
	}
	\haidong{2}
\end{ex}

\begin{ex}
	Phát biểu nào sau đây là đúng khi nói về mô hình động học phân tử?
	\choice
	{Các phân tử cấu tạo nên các chất luôn đứng yên tại một vị trí cố định}
	{\True Các phân tử cấu tạo nên các chất luôn chuyển động không ngừng}
	{Các phân tử cấu tạo nên các chất có thể chuyển động hoặc đứng yên tùy vào đó là chất rắn, chất lỏng hay chất khí}
	{Các phân tử cấu tạo nên chất lỏng và chất khí thì chuyển động, các hạt cấu tạo nên chất rắn thì đứng yên}
	\loigiai{
		\begin{itemize}
			\item Theo mô hình động học phân tử, các phân tử cấu tạo nên các chất luôn chuyển động không ngừng.
			\item Chuyển động của các phân tử xảy ra ngay cả trong trạng thái rắn, lỏng và khí, mặc dù theo những cách khác nhau.
		\end{itemize}
	}
\end{ex}

\begin{ex}
	Tính chất nào sau đây \textbf{không} phải là của phân tử?
	\choice
	{Chuyển động không ngừng}
	{Giữa các phân tử có khoảng cách}
	{\True Có lúc đứng yên, có lúc chuyển động}
	{Chuyển động càng nhanh thì nhiệt độ của vật càng cao}
	\loigiai{
		\begin{itemize}
			\item Các nguyên tử, phân tử chuyển động không ngừng
		\end{itemize}
	}
	\haidong{2}
\end{ex}

\begin{ex}
	Phát biểu nào sau đây \textbf{không} đúng?
	\choice
	{Các chất được cấu tạo từ các hạt gọi là phân tử}
	{Lực tương tác giữa các phân tử ở thể rắn lớn hơn lực tương tác giữa các phân tử ở thể lỏng và thể khí}
	{Các phân tử chất lỏng dao động xung quanh các vị trí cân bằng không cố định}
	{\True Các phân tử đứng sát nhau và giữa chúng không có khoảng cách}
	\loigiai{
		\begin{itemize}
			\item Phát biểu "Các phân tử đứng sát nhau và giữa chúng không có khoảng cách" là không đúng. Trong thực tế, ngay cả khi các phân tử ở gần nhau nhất, giữa chúng vẫn tồn tại khoảng cách nhất định do sự tồn tại của các lực tương tác và các quy luật cơ học lượng tử. 
			\item Những phát biểu còn lại đều đúng:
			\begin{itemize}
				\item "Các chất được cấu tạo từ các hạt gọi là phân tử" là đúng vì theo mô hình cấu tạo của vật chất, tất cả các chất đều được cấu thành từ các phân tử.
				\item "Lực tương tác giữa các phân tử ở thể rắn lớn hơn lực tương tác giữa các phân tử ở thể lỏng và thể khí" là đúng vì ở thể rắn, các phân tử sắp xếp gần nhau và có lực hút lớn, tạo thành cấu trúc cố định, trong khi ở thể lỏng và thể khí, các lực tương tác yếu hơn.
				\item "Các phân tử chất lỏng dao động xung quanh các vị trí cân bằng không cố định" cũng là đúng vì trong chất lỏng, các phân tử có vị trí cân bằng thay đổi liên tục do chúng có thể trượt qua nhau nhưng không tách rời hoàn toàn.
			\end{itemize}
		\end{itemize}
	}
\end{ex}

\begin{ex}
	Phát biểu nào sau đây là \textbf{không} đúng với mô hình động học phân tử?
	
	\choice
	{Các chất được cấu tạo từ các hạt riêng biệt là phân tử}
	{Các phân tử chuyển động không ngừng}
	{\True Tốc độ chuyển động của các phân tử cấu tạo nên vật càng lớn thì thể tích của vật càng lớn}
	{Giữa các phân tử có lực tương tác gọi là lực liên kết phân tử}
	\loigiai{
		\begin{itemize}
			\item Các phân tử chuyển động càng nhanh thì nhiệt độ của vật do chúng tạo nên càng cao
		\end{itemize}
		
	}
	\haidong{2}
\end{ex}

\begin{ex}
	Lực liên kết giữa các phân tử	
	\choice
	{là lực hút}
	{là lực đẩy}
	{tuỳ thuộc vào thể của nó, ở thể rắn là lực hút còn ở thể khí là lực đẩy}
	{\True gồm cả lực hút và lực đẩy}
	\loigiai{
		\begin{itemize}
			\item Giữa các phân tử có lực hút và lực đẩy, gọi chung là lực liên kết phân tử.
		\end{itemize}
	}
	\haidong{2}
\end{ex}

\begin{ex}
	Phát biểu nào sau đây nói về lực tương tác phân tử là \textbf{không} đúng?
	\choice
	{Lực tương tác phân tử đáng kể khi các phân tử ở rất gần nhau}
	{Lực hút phân tử có thể lớn hơn lực đẩy phân tử}
	{\True Lực hút phân tử không thể lớn hơn lực đẩy phân tử}
	{Lực hút phân tử có thể bằng lực đẩy phân tử}
	\loigiai{
		\begin{itemize}
			\item Giữa các phân tử luôn có cả lực hút và lực đẩy. Lực hút và sẽ lớn hơn hay nhỏ hơn lực đẩy tùy vào khoảng cách giữa các phân tử.
		\end{itemize}
	} 
	\haidong{2}
\end{ex}












\begin{ex}
	Đặc điểm nào sau đây là của chất rắn?
	\choice
	{\True Lực tương tác giữa các phân tử rất mạnh}
	{Không có hình dạng xác định}
	{Lực tương tác giữa các phân tử rất yếu}
	{Không có thể tích riêng xác định}
	\loigiai{
		\begin{itemize}
			\item Trong chất rắn, các phân tử ở rất gần nhau nên lực tương tác giữa các phân tử chất rắn rất mạnh.
		\end{itemize}
	}
	\haidong{2}
\end{ex}

\begin{ex}
	Vật rắn có hình dạng xác định vì phân tử cấu tạo nên vật rắn
	\choice
	{không chuyển động}
	{đứng sát nhau}
	{chuyển động với vận tốc nhỏ không đáng kể}
	{\True chuyển động quanh một vị trí cân bằng xác định}
	\loigiai{
		\begin{itemize}
			\item Các phân tử chất rắn được giữ ở các vị trí cân bằng và mỗi phân tử chỉ có thể dao động xung quanh vị trí cân bằng xác định này nên các chất ở thể rắn có thể tích và hình dạng xác định.
		\end{itemize}
	}
	\haidong{2}
\end{ex}

\begin{ex}
	Tại sao người ta \textbf{không} nén chất rắn lại để tiết kiệm diện tích?
	\choice
	{\True Do chất rắn liên kết chặt chẽ với nhau nên rất khó nén}
	{Do chất rắn liên kết không chặt chẽ với nhau nên khó nén}
	{Do các hạt của chất rắn chuyển động tự do}
	{Do người ta không muốn nén chất rắn}
	\loigiai{
		\begin{itemize}
			\item Trong chất rắn, các phân tử ở rất gần nhau nên lực tương tác giữa các phân tử chất rắn chặt chẽ, do đó nó có thể tích xác định và không thể nén được.
		\end{itemize}
	}
	\haidong{2}
\end{ex}

\begin{ex}
	Các nguyên tử trong một miếng sắt có tính chất nào sau đây?
	\choice
	{Khi nhiệt độ tăng thì nở ra}
	{Khi nhiệt độ giảm thì co lại}
	{\True Khoảng cách rất gần nhau}
	{Khoảng cách rất xa nhau}
	\loigiai{
		\begin{itemize}
			\item Các nguyên tử trong kim loại nói riêng và chất rắn nói chung thường đứng rất gần nhau.
		\end{itemize}		
	}
	\haidong{2}
\end{ex}

\begin{ex}
	Các nguyên tử sắp xếp theo một trật tự nhất định trong không gian được gọi là
	\choice
	{mạng nguyên tử}
	{mạng internet}
	{\True mạng tinh thể}
	{mạng nhện}
	\loigiai{
		\begin{itemize}
			\item Các phân tử liên kết chặt với nhau và sắp xếp theo một trật tự hình học xác
			định, tuần hoàn trong không gian, gọi là mạng tinh thể.
		\end{itemize}
	}
	\haidong{2}
\end{ex}

\begin{ex}
	Trong các chất sau, chất nào \textbf{không} phải là chất rắn kết tinh?
	\choice
	{Muối ăn}
	{\True Thuỷ tinh}
	{Kim cương}
	{Thạch anh}
	\loigiai{
		\begin{itemize}
			\item Thủy tinh không có cấu trúc tinh thể, là chất rắn vô định hình
		\end{itemitemize}
	}
	\haidong{2}
\end{ex}

\begin{ex}
	Chất nào sau đây \textbf{không} có cấu trúc tinh thể?
	\choice
	{\True Miếng nhựa thông}
	{Hạt đường}
	{Viên kim cương}
	{Khối thạch anh}
	\loigiai{
		\begin{itemize}
			\item Miếng nhựa thông là chất rắn vô định hình, không có cấu trúc tinh thể
		\end{itemize}
	}	
	\haidong{2}
\end{ex}

\begin{ex}
	Đặc điểm nào sau đây là đặc điểm của thể lỏng?	
	\choice
	{Khoảng cách giữa các phân tử rất lớn so với kích thước của chúng}
	{\True Lực tương tác phân tử yếu hơn lực tương tác phân tử ở thể rắn}
	{Không có thể tích và hình dạng riêng xác định}
	{Các phân tử dao động xung quanh vị trí cân bằng xác định}
	\loigiai{
		\begin{itemize}
			\item Khoảng chất lỏng, các phân tử ở xa nhau hơn so với các phân	tử trong chất rắn nên lực tương tác giữa các phân tử chất lỏng nhỏ hơn trong chất
			rắn
		\end{itemize}
		
	}
	\haidong{3}
\end{ex}

\begin{ex}
	Chất lỏng có thể tích xác định, nhưng hình dạng không xác định là do trong chất lỏng có những đặc điểm nào sau đây?
	\choice
	{Lực liên kết giữa các phân tử chất lỏng là rất lớn, các phân tử chỉ dao động không ngừng quanh một vị trí xác định}
	{Lực liên kết giữa các phân tử chất lỏng là rất yếu, các phân tử dao động tự do về mọi phía}
	{\True Lực liên kết giữa các phân tử trong chất lỏng yếu hơn so với chất rắn, các phân tử khi dao động còn có thể trượt lên nhau}
	{Lực liên kết giữa các phân tử chất lỏng là rất yếu, các phân tử dao động xung quanh một vị trí xác định}
	\loigiai{
		\begin{itemize}
			\item Lực liên kết giữa các phân tử yếu hơn trong chất rắn nên các chất ở thể lỏng có thể tích xác định nhưng không có hình dạng riêng mà có hình dạng
			của phần bình chứa nó
		\end{itemize}
	}
	\haidong{1}
\end{ex}















\begin{ex}
	Đặc điểm nào sau đây là đặc điểm của chất lỏng?
	\choice
	{Có hình dạng của phần bình chứa nó, có thể rót được và chảy tràn trên bề mặt, dễ nén}
	{\True Có hình dạng của phần bình chứa nó, có thể rót được và chảy tràn trên bề mặt, rất khó nén}
	{Có hình dạng cố định, có thể rót được và chảy tràn trên bề mặt, rất khó nén}
	{Có hình dạng của phần bình chứa nó, dễ dàng lan tỏa trong không gian, dễ nén}
	\loigiai{
		\begin{itemize}
			\item Do lực liên kết giữa các phân tử yếu hơn trong chất rắn nên các chất ở thể lỏng có thể tích xác định nhưng không có hình dạng riêng mà có hình dạng của phần bình chứa nó.
		\end{itemize}
	}
	\haidong{2}
\end{ex}

\begin{ex} \nguon{1}
	\immini[thm]{Chuyển động hỗn loạn, không ngừng của các hạt rất nhỏ dưới tác dụng của các phân tử chất lỏng hoặc chất khí được gọi là chuyển động Brown}
	{\includegraphics[width=3cm]{images/2024_05_26_010fc85c05a422f9b1ffg-01(1)}}
	\choiceTF[t]
	{Chuyển động Brown do một nhà Vật Lí học phát hiện ra}
	{Khi quan sát các hạt phấn hoa, Brown khẳng định cấu tạo chất bao gồm các phân tử}
	{\True Nguyên nhân của chuyển động Brown là do các phân tử nước luôn luôn chuyển động hỗn loạn, va chạm vào các hạt phấn hoa từ nhiều phía}
	{Hạt phấn hoa chuyển động có quy luật xác định}	
	\loigiai{
		\begin{itemchoice}
			\itemch 
			\itemch 
			\itemch 
			\itemch 
		\end{itemchoice} 
		Chuyển động Brown được phát hiện khi quan sát các hạt phấn hoa chuyển động hỗn loạn không có quy luật ở trong nước.
		Nguyên nhân là do các phân tử nước luôn luôn chuyển động hỗn loạn, va chạm vào các hạt phấn hoa từ nhiều phía.
	}
\end{ex}

\begin{ex} \nguon{4}
	Mô hình động học phân tử có thể dùng để giải thích cấu trúc của các chất rắn, chất lỏng và chất khí. Các chất được cấu tạo từ các nguyên tử, phân tử và chúng chuyển động không ngừng. Sự sắp xếp và tương tác khác nhau giữa các phân tử sẽ tạo nên các tính chất khác nhau cho chất rắn, chất lỏng, chất khí.
	\choiceTF[t]
	{\True giữa các phân tử có lực tương tác. Độ lớn của những lực này phụ thuộc vào khoảng cách giữa chúng}
	{\True các phân tử chuyển động không ngừng, chuyển động này càng nhanh thì nhiệt độ của vật càng cao}
	{\True các chất được cấu tạo từ những hạt có kích thước rất nhỏ gọi là phân tử. Giữa các phân tử có khoảng cách}
	{khi khoảng cách giữa các phân tử lớn thì lực đẩy mạnh hơn lực hút}
	\loigiai{
		\begin{itemchoice}
			\itemch 
			\itemch 
			\itemch 
			\itemch Khi khoảng cách giữa các phân tử lớn thì lực hút mạnh hơn lực đẩy.
		\end{itemchoice}
	}
	\haidong{10}
\end{ex}


\begin{ex} \nguon{4}
	Mô hình động học phân tử có thể dùng để giải thích cấu trúc của các chất rắn, chất lỏng và chất khí. Các chất được cấu tạo từ các nguyên tử, phân tử và chúng chuyển động không ngừng. Sự sắp xếp và tương tác khác nhau giữa các phân tử sẽ tạo nên các tính chất khác nhau cho chất rắn, chất lỏng, chất khí.
	\choiceTF[t]
	{\True giữa các phân tử có lực tương tác. Độ lớn của những lực này phụ thuộc vào khoảng cách giữa chúng}
	{\True các phân tử chuyển động không ngừng, chuyển động này càng nhanh thì nhiệt độ của vật càng cao}
	{\True các chất được cấu tạo từ những hạt có kích thước rất nhỏ gọi là phân tử. Giữa các phân tử có khoảng cách}
	{khi khoảng cách giữa các phân tử lớn thì lực đẩy mạnh hơn lực hút}
	\loigiai{
		\begin{itemchoice}
			\itemch 
			\itemch 
			\itemch 
			\itemch Khi khoảng cách giữa các phân tử lớn thì lực hút mạnh hơn lực đẩy.
		\end{itemchoice}
	}
	\haidong{10}
\end{ex}


\begin{ex} \nguon{4}
	Xét cấu trúc của chất khí thì
	\choiceTF[t]
	{\True khoảng cách giữa các phân tử chất khí rất lớn so với kích thước của chúng}
	{\True trừ khi va chạm, lực tương tác giữa các phân tử khí rất nhỏ, hầu như không đáng kể}
	{một lượng không khí xác định luôn có thể tích và hình dạng riêng xác định}
	{\True các phân tử chất khí chuyển động hỗn loạn, không ngừng về mọi phía, chiếm toàn bộ không gian bình chứa}
	\loigiai{
		\begin{itemchoice}
			\itemch 
			\itemch 
			\itemch Chất khí không có thể tích và hình dạng xác định
			\itemch 
		\end{itemchoice}
	}
	\haidong{5}
\end{ex}


\begin{ex} \nguon{4}
	Kim cương và than chì đều có cấu trúc dưới dạng chất rắn, cùng được cấu tạo từ nguyên tố carbon nhưng lại có tính chất khác nhau, điển hình kim cương rất cứng còn than chì thì mềm. 
	\choiceTF[t]
	{\True Sự khác nhau về tính chất của kim cương và than chì là do cấu trúc mạng tinh thể của chúng}
	{\True Mỗi nguyên tử carbon trong tinh thể kim cương liên kết với bốn nguyên tử khác tạo nên một cấu trúc tinh thể vững chắc}
	{\True Than chì có cấu trúc tinh thể dạng xếp thành lớp, mỗi nguyên tử carbon liên kết với ba nguyên tử carbon khác. Cấu trúc này không tạo nên sự bền vững cho than chì}
	{Do kim cương là chất rắn kết tinh còn than chì là chất rắn vô định hình}
	\loigiai{
		\begin{itemchoice}
			\itemch 
			\itemch 
			\itemch 
			\itemch 
		\end{itemchoice}
	}
\end{ex}